\documentclass[9pt,letterpaper,onecolumn]{rho-class/rho}
%\usepackage[english]{babel}
\usepackage[spanish,es-nodecimaldot,es-noindentfirst]{babel}

\setbool{rho-abstract}{true} % Set false to hide the abstract
\setbool{corres-info}{true} % Set false to hide the corresponding author section

%----------------------------------------------------------
% TITULO
%----------------------------------------------------------

\title{Proyecto 2: Divide y venceras}

%----------------------------------------------------------
% AUTORES
%----------------------------------------------------------

\author[$\dagger$]{Felipe Paillacar}
\author[$\dagger$]{Constanza Araya}
\author[$\dagger$]{Guillermo Vargas}

%----------------------------------------------------------

%\affil[1]{Affiliation of author one}
%\affil[2]{Affiliation of author two}
%\affil[3]{Affiliation of author three}
\affil[$\dagger$]{Universidad de Magallanes}
\dates{\today}
\smalltitle{Diseño de algoritmos}
\theday{\today}
%----------------------------------------------------------
% ABSTRACT
%----------------------------------------------------------

\begin{abstract}
    Este proyecto aborda la implementación y el análisis experimental de algoritmos basados en la estrategia de divide y vencerás
    integrados a un sistema de gestión de deportistas. El estudio se centra en evaluar el rendimiento de técnicas de ordenamiento 
    como Merge Sort (con optimización de umbral) y Quick Sort (con diversas variantes de pivote), junto con métodos de búsqueda binaria y
    exponencial. El trabajo integra un análisis de complejidad teórica contrastado con pruebas empíricas de tiempo de ejecución para 
    determinar la eficiencia de cada algoritmo.
\end{abstract}

%----------------------------------------------------------

%\keywords{keyword 1, keyword 2, keyword 3, keyword 4, keyword 5}

%----------------------------------------------------------

\begin{document}
	
    \maketitle
    \thispagestyle{firststyle}
    \tableofcontents
    %\linenumbers
\newpage
%----------------------------------------------------------
\section{Objetivos}
\begin{itemize}
    \item Implementar correctamente algoritmos utilizando la estrategia de divide y venceras.
    \item Analizar empiricamente la complejidad y el rendimiento de los algoritmos implementados, evaluando ademas el impacto de sus distintas optimizaciones y variantes.
    \item Documentar adecuadamente el codigo desarrollado y los resultados obtenidos del analisis.
\end{itemize}

\newpage
%-----------------------------------------------------------------------------------------------------------------------------------%
\section{Introduccion}




\newpage
%-----------------------------------------------------------------------------------------------------------------------------------%
\section{Conceptos teoricos}
    \subsection{Divide y venceras}
    Divide y venceras es una estrategia algoritmica. Un algoritmo divide y venceras tipico resuelve un problema siguiendo tres pasos:
        \begin{itemize}
            \item \textbf{Dividir:} Descompones el problema en subproblemas mas pequeños, cada subproblema debe representar una parte del problema
            original, Por lo general este paso emplea un enfoque recursivo para dividir el problema hasta que no es posible crear subproblemas mas pequeños.
            \item \textbf{Vencer:} Resolver los sub-problemas recursivamente. Esto se conoce como el caso base, es decir el subproblema es tan pequeño 
            que se 'resolvio solo'.
            \item \textbf{Combinar:} Combinar apropiadamente las soluciones de los subproblemas para obtener la solucion del problema original.
        \end{itemize}

    \subsection{Complejidad de un algorimo}
    La complejidad es una medida de que tan eficiente es el algoritmo para resolver un problema, es decir es una medida de cuanto tiempo y espacio
    (memoria) requiere el algoritmo para producir una solucion.
    La complejidad de un algoritmo se expresa utilzando la notacion \textbf{Big O}. Esta notacion proporciona una forma de comparar la complejidad
    de distintos algoritmos. Por ejemplo, un algoritmo con una complejidad temporal de $O(n)$ se considera más eficiente que un algoritmo con 
    una complejidad temporal de $O(n^2)$, porque crece a un ritmo más lento a medida que aumenta el tamaño de entrada.
    En general, el objetivo del diseño de algoritmos es encontrar algoritmos que sean eficientes y fáciles de implementar. 
    Esto implica un compromiso entre la complejidad del tiempo y el espacio, así como otros factores como la simplicidad (facil entendimiento) 
    y la mantenibilidad (facilidad de actualización, corrección o mejora).

        \subsubsection{Conceptos clave de complejidad (notacion Big O)}
        Es comun confundir la notacion Big O como una medida de tiempo, pero la realidad es que la notacion Big O nos indica cuantas operaciones
        debe realizar un algoritmo para resolver un problema, ya que si fuera una medida de tiempo podria variar dependiendo del hardware en el que se estan
        realizando las pruebas. Por ello se habla del mejor y peor caso, es decir el mejor caso o caso ideal seria cuando un algoritmo tiene que hacer operaciones minimas
        para llegar a la solucion. Por ejemplo, en el caso de la busqueda lineal, el mejor caso es cuando el elemento que se busca esta en la primera posicion del arreglo, 
        por lo que solo se realiza una operacion. A continuacion se muestra una comparacion de las funciones de complejidad mas comunes.
            \begin{figure}[H]
                \centering
                \includegraphics[width=0.4\columnwidth]{caption1.png}
                \caption{Comparacion de las funciones de complejidad mas comunes.}
                \label{fig:comparacion-complejidad}
            \end{figure}
            
            \begin{itemize}
                \item \textbf{$O(1):$} Tiempo constante. El tiempo de ejecución no cambia independiente del tamaño de la entrada.
                \item \textbf{$O(\log n):$} Tiempo logaritmico, muy eficiente. El tiempo de ejecución crece lentamente a medida que aumenta el 
                tamaño de la entrada.
                \item \textbf{$O(n):$} Tiempo lineal. El tiempo de ejecución crece proporcionalmente al tamaño de la entrada.
                \item \textbf{$O(n \log n):$} Tiempo log-lineal. Es casi una línea recta, pero con una ligera inclinación hacia arriba.
                \item \textbf{$O(n^2):$} Tiempo cuadratico. El tiempo de ejecución crece proporcionalmente al cuadrado del tamaño de la entrada.
            \end{itemize}

    \subsection{Esquema de particion de Lomuto} \label{ssec:lomuto}
    El esquema de particion de Lomuto es un algoritmo popular utilizado para dividir una matriz. Lomuto funciona seleccionando un elemento
    de pivote, generalmente el ultimo elemento del array, y luego dividiendolo en dos subarrays de modo que los elementos menores o iguales
    al pivote esten en el lado izqueirdo del pivote y los elementos mayores queden a su derecha. El esquema de particion de Lomuto
    funciona de la siguiente manera:
        \begin{itemize}
            \item Se selecciona un elemento del array como pivote (generalmente el ultimo).
            \item Se recorre la lista comparando cada elemento con el pivote. Los menores se mueven a la izquierda y los mayores a la derecha.
            \item Al final del recorrido, se coloca el pivote en su lugar definitivo (entre las dos sublistas)
            \item Se repite el mismo proceso en la sublista izquierda y en la derecha hasta que todos los elementos estén ordenados.
        \end{itemize}
            

\section{Algoritmos de ordenamiento}
    Los algoritmos de ordenamiento son un conjunto de procedimientos que organizan los elementos de una lista o arreglo en un orden 
    específico, ya sea ascendente o descendente.
    Los algoritmos de ordenamiento son importantes ya que facilitan la busqueda, la organizacion y el analisis de datos. Son fundamentales
    en la eficiencia informatica ya que el desorden es costoso en terminos de tiempo y recursos. Algunas razones de su importancia incluyen:
    \begin{itemize}
        \item \textbf{Optimizacion de la busqueda}
            \begin{itemize}
                \item \textbf{Busqueda lineal:} Si los datos estan desordenados, se debe revisar uno por uno hasta encontrar el elemento deseado (complejidad $O(n)$). 
                En cambio, si los datos estan ordenados, se pueden utilizar algoritmos de busqueda mas eficientes.
                \item \textbf{Busqueda binaria:} Este algoritmo requiere que los datos esten ordenados para funcionar correctamente. 
                Permite encontrar un elemento en tiempo logaritmico $O(\log n)$.
            \end{itemize}
        \item \textbf{Organizacion y visualizacion:} El ordenamiento permite agrupar elementos similares, lo cual es vital para:
            \begin{itemize}
                \item Eliminar duplicados: Es mas facil notar que un elemento esta duplicado si esta junto a su par.
                \item Priorizacion: En sistemas operativos o servidores, ordenar procesos por prioridad asegura que lo mas importante se ejecute primero.
            \end{itemize} 
    \end{itemize}

    \begin{info}
        Antes de elegir un algoritmo de ordenamiento es importante hacer algunas preguntas: ¿Cuán grande son los datos que se ordenan?
        ¿Cuánta memoria hay disponible? ¿Los datos creceran?
    \end{info}
        
    \subsection{Merge Sort}\label{ssec:merge-sort}
        \begin{figure}[H]
            \centering
            \includegraphics[width=0.4\columnwidth]{caption2.png}
            \caption{Merge sort no adaptativo}
            \label{fig:merge-sort}
        \end{figure}

        Merge sort es un algoritmo de ordenamiento conocido por su eficiencia y estabilidad. Utiliza la estrategia de divide y venceras 
        para ordenar una lista o arreglo. Funciona dividiendo recursivamente el arreglo en mitades hasta que cada subarreglo contiene 
        un solo elemento, luego combina esos subarreglos de manera ordenada para formar el arreglo final ordenado.
        Merge sort es un algoritmo no adaptativo, lo que significa que su rendimiento no mejora si el array esta previamente ordenado. 
        Debido a esto su complejidad en el peor y mejor caso es $O(n \log n)$

    \subsection{Merge Sort con Insertion Sort}\label{ssec:merge-sort-insertion}
        Como se menciono anteriormente en \ref{ssec:merge-sort}, merge sort es un algoritmo no adaptativo, lo que significa que su 
        rendimiento no mejora si el array esta previamente ordenado. Para ello se puede optimizar utilizando insertion sort para ordenar
        subarreglos pequeños, ya que insertion sort es un algoritmo adaptativo, lo que significa que su rendimiento mejora si el array 
        esta previamente ordenado, es decir Insertion sort no es 'ciego' como merge sort
        Esto ayuda a que en teoria en el peor caso se mantenga la complejidad de merge sort $O(n \log n)$, pero en el mejor caso se 
        mejora a $O(n)$, ya que insertion sort detecta que los subarreglos estan ordenados y como los subarreglos son pequeños, mejora 
        significativamente el tiempo de ejecucion.


    \subsection{Quick Sort}
        \begin{figure}[H]
            \centering
            \includegraphics[width=0.4\columnwidth]{caption3.png}
            \caption{Quick sort}
            \label{fig:quick-sort}
        \end{figure}
        Quick sort es un algoritmo de ordenamiento que utiliza la estrategia de divide y venceras que elige un elemento como pivote
        y divide el array al rededor del pivote, como se nombro en \ref{ssec:lomuto}. Este algoritmo utiliza recursivudad, para ello es 
        importante tener claro cual es el caso base, el cual es cuando el array tiene un solo elemento o esta vacio, 
        ya que en ese caso el array ya esta ordenado.
        Hay diferentes formas de elegir un pivote, como por ejemplo:
        \begin{itemize}
            \item \textbf{Pivote primer elemento}
            \item \textbf{Pivote ultimo elemento}
            \item \textbf{Pivote aleatorio}
            \item \textbf{Pivote mediana de tres}
        \end{itemize}
        La eleccion del pivote es importante ya que puede afectar significativamente el rendimiento del algoritmo. Lo cual sera analizado
        en el transcurso de este proyecto.


\newpage
%-----------------------------------------------------------------------------------------------------------------------------------%

\section{Algoritmos de busqueda}








\newpage
%-----------------------------------------------------------------------------------------------------------------------------------%

\section{Algoritmos de seleccion}











\newpage
%-----------------------------------------------------------------------------------------------------------------------------------%
\section{Analisis de casos: ordenamiento}














\newpage
%-----------------------------------------------------------------------------------------------------------------------------------%
\section{Analisis de casos: busqueda}











\newpage
%-----------------------------------------------------------------------------------------------------------------------------------%
\section{Analisis de casos: seleccion}






\newpage
%-----------------------------------------------------------------------------------------------------------------------------------%
\section{Conclusion}




\begin{lstlisting}[language=C, caption=Here is an implementation of a delete function for a singly linked list in C.(\href{https://users.cs.fiu.edu/~weiss/\#dsaac2e}{Source code})]
    /* Delete from a list */
/* Cell pointed to by P->Next is wiped out */
    /* Assume that the position is legal */
    /* Assume use of a header node */

    void
    Delete( ElementType X, List L )
    {
        Position P, TmpCell;

        P = FindPrevious( X, L );

        if( !IsLast( P, L ) )  /* Assumption of header use */
        {                      /* X is found; delete it */
            TmpCell = P->Next;
            P->Next = TmpCell->Next;  /* Bypass deleted cell */
            free( TmpCell );
        }
    }
\end{lstlisting}


%----------------------------------------------------------
\newpage
\nocite{*}
\printbibliography

%----------------------------------------------------------

\end{document}